% !TEX root = ../bb_AiRa_Kappaleet.tex

% ö = \%c3\%b6
% ä = \%C3\%A4

\xdef\sectiontitle{\IVsectiontitle} %aiheuttama
\TOCrefs

%%%%%%%%%%%%%%%
\begin{frame}[label=4_4_a]
\frametitle{\texorpdfstring{\culine[\sectiontitleulinecolor]{\sectiontitle}}{\sectiontitle} }
\label{\TOCname}

\vspace{-3mm}

%\vspace{\chapterTOCvksip}%
\begin{itemize}[leftmargin=0.5cm,itemsep=\chapterTOCitemsep]
    \item[4.1]<1-> \hyperlink{4_1_a}{\IVaSubsectiontitle}
    \item[4.2]<1-> \hyperlink{4_2_a}{\IVbSubsectiontitle}	
    \item[4.3]<1-> \hyperlink{4_3_a}{\IVcSubsectiontitle}	
    \item[4.4]<1-> \hyperlink{4_4_a}{\CDAlert[red]{\IVdSubsectiontitle}}
    \item[4.5]<1-> \hyperlink{4_5_a}{\IVeSubsectiontitle}
    \item[4.6]<1-> \hyperlink{4_6_a}{\IVfSubsectiontitle}
    \item[4.7]<1-> \hyperlink{4_7_a}{\IVgSubsectiontitle}
\end{itemize}


\end{frame}
%%%%%%%%%%%%%%%

\xdef\sectiontitle{\IVdSubsectiontitle} 

%%%%%%%%%%%%%%%
\begin{frame}
\frametitle{\texorpdfstring{\culine[\sectiontitleulinecolor]{\sectiontitle}}{\sectiontitle} }
\framesubtitle{Lyhyt historia}%A Short history
\appear{}

\semismall
\begin{itemize}
\item \appear{Varhaiset kokeet tehtiin \CDAlert[red]{alfalähteillä}} \appear{esim.\ Pb-210}\appear{, \CDAlert[red]{5,3\;MeV:n kineettisillä energioilla}:} 

\item<.->[] \begin{itemize}
	 \item \CDAlert[red]{Elektroni}\appear{, löydettiin 1897} \appear{((J. J. Thomson, katodisäteellä, e/m))}\appear{, sen varaus mitattiin 1909 (Millikan)}
	\item \CDAlert[blue]{Protoni}\appear{, löydettiin 1919} \appear{((E. Rutherford, alfahiukkaspommituksella kultakalvoon))}
	\item \CDAlert[fuchsia]{Neutroni}\appear{, löydettiin 1932} \appear{((J. Chadwick, sähköisesti neutraalien painavien hiukkasten säteily löydettiin pommituskokeissa, massa oli lähellä protonia))} \appear{$\tau \sim 882 \mathrm{\;s}$}
\end{itemize}

\item Myöhemmin \CDAlert[blue]{gammasäteitä ja kosmisia säteitä} käytettiin $\appear{\oldrightarrow$  \CDAlert[blue]{$\gtrsim 1 \mathrm{\;GeV}$}:}
\item<.->[] \begin{itemize}
\item \CDAlert[red]{Positroni}\appear{, löydettiin 1932} \appear{((Anderson, Wilsonin sumukammiota \& gammasäteitä))}
\item Myoni\appear{, löydettiin 1936 (kosmisilla säteillä)} \appear{$m_{\mu}=0,106 \mathrm{\;GeV}, ~\tau \sim 2,2 \;\mu\mathrm{s}$}
\item Pioni\appear{, löydettiin 1947}\appear{, (kosmisilla säteillä) $m_{\pi}=0,140 \mathrm{\;GeV}, ~\tau \sim 26 \mathrm{\;ns}$}
\item \appear{Kaoni}\appear{, löydettiin 1947}\appear{, (kosmisilla säteillä) $m_{K}=0,490 \mathrm{\;GeV}, ~\tau \sim 5 \mathrm{\;ns}$}
\end{itemize}

\item Kaonit $\oldrightarrow$ koostuvat \CDAlert[fuchsia]{kvarkeista}: \appear{ylös, alas, \& outo vuonna 1967}
\end{itemize}

\endofslide 
%\endofsection
\end{frame}
%%%%%%%%%%%%%%%

%%%%%%%%%%%%%%%
\begin{frame}
\frametitle{\texorpdfstring{\culine[\sectiontitleulinecolor]{\sectiontitle}}{\sectiontitle} }
\framesubtitle{Modernit kiihdyttimet}%Modern accelerators
\appear{}

\semismall
\begin{itemize}
	\item Vielä myöhemmin, \appear{uusia ja painavampia hiukkasia löydettiin:}
\begin{itemize} \small
	 \item 1974: \appear{$3,1 \mathrm{\;GeV} ~J/\psi$-\CDAlert[fuchsia]{mesoni}}\appear{, koostuu lumo-kvarkista ja sen antikvarkista}
	\item 1975: \appear{$1,8 \mathrm{\;GeV}$ tauoni}
	\item 1977: \appear{$9,5 \mathrm{\;GeV}$ Upsilon-mesoni}\appear{, koostuu alas-kvarkista ja sen antikvarkista}
	\item 1983: \appear{$80 \mathrm{\;GeV} ~\&~ 91 \mathrm{\;GeV}$ $W$- ja $Z$-bosonit}
	\item 1994: \appear{$180 \mathrm{\;GeV}$ \CDAlert[fuchsia]{huippu-kvarkki}}\appear{, massiivisin kvarkki mikä \CDAlert[fuchsia]{ei 'hadronisoidu'}}
\end{itemize}

%\item Modernit kiihdyttimet yltävät >500\;TeV\appear{, joko käyttäen \CDAlert[blue]{2-säteen kiihdytin}menetelmää} \appear{tai \CDAlert[red]{paikallaan olevaa kohdetta ja yhtä kiihdytettyä sädettä}}
\item Modernit kiihdyttimet yltävät jopa 6,8 TeV:iin ja törmäyttävät kaksi korkeaenergistä hiukkassuihkua toisiinsa ($\Rightarrow$ törmäysenergia jopa 13,6 TeV) 
\end{itemize}

%
\appear{\xdef\loc{south} %south west
\begin{tikzpicture}[remember picture,overlay] % anchor=south
    \node (A) [yshift=-0.15*\figYshift,xshift=-\figXshift, anchor=\loc] at (current page.\loc) {\includegraphics[width=6.6cm]{Figures_booklet/SOM_11_Fundamental_particles_figures/SOM_Fundamental_particles_Figure_4.png}}; %scale=\figscale. % 8cm max hei
\end{tikzpicture}}

%\xdef\loc{south east}
%\begin{tikzpicture}[remember picture,overlay] % anchor=south
%    \node (A) [yshift=-0.15*\figYshift,xshift=\figXshift, anchor=\loc] at (current page.\loc) {\includegraphics[width=5.6cm]{Figures_booklet/SOM_11_Fundamental_particles_figures/SOM_Fundamental_particles_Figure_6.png}}; %scale=\figscale. % 8cm max hei
%\end{tikzpicture}\CR[ref=h]{}}

\endofslide 
%\endofsection
\end{frame}
%%%%%%%%%%%%%%%

%%%%%%%%%%%%%%%
\begin{frame}
\frametitle{\texorpdfstring{\culine[\sectiontitleulinecolor]{\sectiontitle}}{\sectiontitle} }
\framesubtitle{Hiukkaskiihdyttimet}%Particle detectors
\appear{}

\seminormal
\begin{itemize}
	 \item Sumukammioiden jäädessä historiaan\appear{, on useita tekniikoita tunnistaa} \appear{ja karakterisoida törmäyksissä syntyviä/vapautuvia hiukkasia}\appear{, varauksen määrittäminen}\appear{, liikemäärä}\appear{, massa jne.}

\item \CDAlert[fuchsia]{Tuikeilmaisin}: \appear{materiaali emittoi valoa hiukkasen osuessa (esim.\  Th:NaI kide)}
\item \CDAlert[black]{Monijohdin verrannollisuuslaskuri}: \appear{saapuvat hiukkaset ionisoivat ainetta, ja johdinlankaristikoista mitattujen sähkövirtojen avulla hiukkasen liikerata voidaan selvittää, \CDAlert[black]{fysiikan Nobelin palkinto 1992}}
\item \CDAlert[red]{Cerenkov-ilmaisin}: \appear{hiukkasen nopeus voidaan määrittää syntyneen säteilyn ominaisuuksista}
\item \CDAlert[blue]{Pii-ilmaisin}: \appear{Säteily synnyttää puolijohteeseen elektroni–aukkopareja} \appear{$\oldrightarrow$ hiukkasen liikerata}
\item \CDAlert[green]{Kalorimetri}: \appear{Tulevan hiukkasen synnyttämien hiukkassuihkujen avulla selvitetään energia}

\end{itemize}

%\endofslide 
\endofsection
\end{frame}
%%%%%%%%%%%%%%%
